\section{Variables aleatorias discretas}
Para la deducción de las definiciones y ecuaciones de esta sección se seguirá
según~\cite{swendsen2012introduction} en el capítulo 3.

Un evento elemental es un resultado individual y exlusivo de un experimento, tal
que puede ocurrir una y solamente uno a la vez, Ejem. cuando se lanza una
moneda se tiene el resultado de que la moneda salga cara o sello.

\subsection{Variable aleatorias}
Sea $A = \set{a_j}_{j = 1}^N$ un conjunto de eventos elementales donde a cada
evento $a_j$ se le asigna una probabilidad $P(a_j)$, que por
simplicidad escribiremos (A, P(a)), a esto llamaremos
\textit{Variable aleatoria discreta}, si los eventos elementales son
números entonces a $(A, P(a))$ se le llama \textit{Número aleatorio}.
Por definicion de probablidad, para cada elemento se debe cumplir

\begin{equation}
  0 \leq P(a_j) \leq 1
  \label{eq:prob-elemental}
\end{equation}

Asi la probabilidad de todo el conjunto $A$ debe ser $1$, por lo que
debe, cumplirse
que

\begin{equation}
  \sum_{j = 1}^N P(a_j) = 1
  \label{eq:prob-total}
\end{equation}

Que por simplicidad escribiremos

\begin{equation}
  \sum_{a} P(a) = 1
  \label{eq:prob-total-sim}
\end{equation}

\subsection{Probablidad de multiples variables aleatorias discretas}
Supongamos que se tienen dos eventos elementales
$A=\set{a_i}_{i=1}^N$ y $B=\set{b_j}_{j=1}^M$,
entonces la probabilidad conjunta para $a_i$ y $b_j$ es

\begin{equation}
  0 \leq P(a_i, b_j) \leq 1
  \label{eq:prob-conjunta}
\end{equation}

y la probabilidad para todo el cunjunto de eventos elementales se debe cumplir

\begin{equation}
  \sum_{i=1}^N \sum_{j=1}^M P(a_i, b_j) = 1
  \label{eq:prob-conjunta-total}
\end{equation}

que por simplicidad escribiremos

\begin{equation}
  \sum_{a} \sum_{b} P(a, b) = 1
  \label{eq:prob-conjunta-total-sim}
\end{equation}

\subsubsection*{Probabilidad marginal}
Si queremos recuperar la probablidad $P_A(a_j)$ al cual denominaremos
\textit{Probablidad marginal}, es lógico que sumaremos todas las
probablidades sobre $B$ así obtenemos

\begin{equation}
  P_A(a) = \sum_{b} P(a, b)
  \label{eq:prob-marginal}
\end{equation}

\subsubsection*{Probabilidad condicional}
Supongamos que queremos la probablidad de que ocurra el evento
$a$ dado que ocurrio $b$ el cual notaremos como $P(a|b)$, asumiendo
que que la probablidad $P_B(b) \neq 0$ entonces esta será

\begin{equation}
  P(a|b) = \frac{P(a,b)}{P_B(b)}
  \label{eq:prob-condicional}
\end{equation}

% === Funciones aleatorias y Números Aleatorios ===
\subsection{Funciones aleatorias de una variable y Números Aleatorios}
Sea el conjunto de eventos elementales $A = \set{a_j}_{j = 1}^{N_A}$ junto a
su probablidad $P(a_i)$, podemos definir sobre este número aleatorio
una función númerica $F = \set{F(a_j)}_{j = 1}^{N_A}$ tal que

\begin{equation}
  P_F(f) = \sum_a \delta_{f,F(a)} P(a)
  \label{eq:f-prob}
\end{equation}

el cual también será un número aleatorio.

Supongamos ahora que se tiene dos números aleatorios $X, Y$, con probablidad
$P(x, y)$ y se quiere una nueva función arbitraria $G(x, y)$, entonces
la distribución de probablidad de $G$ esta dado como una suma sobre
todas las combinaciónes de los eventeos $X$ y $Y$ con la delta de kronecker
tomando cada uno de los eventos particulares correspondientes a $G$

\begin{equation}
  P_G(g) = \sum_{x}\sum_y\delta_{g, G(x, y)} P(x, y)
  \label{eq:prob-g}
\end{equation}

\subsection{Media, Variancia y Desviación Estándar}

La media o promedio de una función $F(A)$, sobre $(A, P(a))$ se define
como

\begin{equation}
  \bk{F} = \sum_a F(a) P(a)
  \label{eq:d-mean}
\end{equation}

De manera similar el $n$-ésimo momento se define

\begin{equation}
  \bk{F^n} = \sum_a F^n(a) P(a)
  \label{eq:d-n-moment}
\end{equation}

Y el $n$-ésimo momento central

\begin{equation}
  \bk{(F - \bk{F})^n}
  = \sum_a (F^n(a) - \bk{F}^n)(a) P(a)
  \label{eq:d-central-moment}
\end{equation}

El segundo momento central se le denomina \textit{varianza}, $\sigma^2$ el cuál
será

\begin{align*}
  \sigma^2 &= \bk{(F - \bk{F})^2}\\
  &= \bk{F^2 - 2F\bk{F} + \bk{F}^2}\\
  &= \sum_a (F^2 - 2F\bk{F} + \bk{F}^2)(a) P(a)\\
  &= \sum_a F^2(a)P(a) - 2\bk{F}\sum_a F(a)P(a) + \bk{F}^2\sum_aP(a)\\
  &= \bk{F^2} - 2\bk{F}\bk{F} + \bk{F}^2
\end{align*}

\begin{equation}
  \therefore \sigma^2 = \bk{F^2} - \bk{F}^2
  \label{eq:d-var}
\end{equation}

Con esto podemos definir la desviación estándar como $\sigma = \sqrt{\sigma^2}$

\begin{equation}
  \sigma = \sqrt{\bk{F^2} - \bk{F}^2}
  \label{eq:d-std}
\end{equation}
